\chapter{Introduction}\label{introduction}
% The introduction of your bachelor thesis introduces the research area, the
% research hypothesis, and the scientific contributions of your work.
% A good narrative structure is the one suggested by Simon Peyton Jones
% \cite{peys04:HowToWriteAGoodResearchPaper}:

% \begin{itemize}
% \item describe the problem / research question
% \item motivate why this problem must be solved
% \item demonstrate that a (new) solution is needed
% \item explain the intuition behind your solution
% \item motivate why / how your solution solves the problem (this is technical)
% \item explain how it compares with related work
% \end{itemize}

% Close the introduction with a paragraph in which the content of the next chapters
% is briefly mentioned (one sentence per chapter). 

% Starting a new paragraph is done by inserting an empty line like this.

One of the simplest, yet most powerful, models of computation is the \lc.  
It serves as the foundation of many programming languages and proof assistants, such as Haskell and Agda.
The two main well-known evaluation strategies for the \lc\ are \emph{call-by-name} (cbn) and \emph{call-by-value} (cbv). 
Call-by-name does not evaluate an argument to a value until it is needed, while call-by-value evaluates an argument to a value before passing it to a function.
These are two very different approaches and depending on the context, one may be more suitable than the other.

Cbn is more efficient when some arguments are not used inside a function, as it avoids unnecessary computation.
However, if an argument is used multiple times, cbv is preferable, as it avoids recomputing the argument each time it is used.
Therefore, a \lc that can express both cbn and cbv is ideal. 
It allows for the creation of a programming language or proof assistant that is more flexible in evaluation strategies.
While cbv can be simulated in Haskell, it requires some workarounds and support for cbv is not built into the language.  

Many approaches of unifying cbn and cbv have been proposed. 
Thomas Ehrhard and Giulio Guerrieri \cite{EhrhardGuerrieri2016} proposed a linear logic approach,
just like the one by Maraist et al. \cite{MARAIST1999175},
while Hatcliff, John and Danny, Olivier \cite{HATCLIFF_DANVY_1997} proposed a thunk-based approach.
However, we focus on the modal approach, introduced by Espírito Santo et al. \cite{espritosanto_et_al:LIPIcs:2019:10525}.
This approach is based on the modal \lc \labox, which is a fragment of Bierman and de Paiva's intuitionistic modal logic (\textsf{IS4}) \cite{BiermanPaiva2000}.
The proposed call-by-box \lc (\lab) is a refinement of \labox, which has simpler typing rules and a more straightforward reduction strategy.
This thesis aims to formalise this approach in Agda, which is a dependently typed programming language and proof assistant.
Although \lab is discussed in the literature, it has not yet been formalised in a proof assistant before.
Given the extra complexity of the \lab in comparison to the standard \lc, 
it is important to formalise it to ensure that the properties of \lab are correct and machine-checked.
Therefore, the main research question of this thesis is as follows:
\begin{quote}
  How can the unification of call-by-name and call-by-value using modal logic be formalised in Agda?
\end{quote} 

In \lab, cbn and cbv can be simulated by using \emph{Girard's} and \emph{Gödel's} embeddings, respectively.
These are embeddings from the standard \lc to \lab, preserving both reduction and evaluation. 
This thesis formalises both embeddings in Agda and we prove the following:

\begin{itemize}
  \item Girard's embedding preserves reduction.
  \item Girard's embedding preserves evaluation.
  \item Gödel's embedding preserves reduction.
\end{itemize}

We did not formally prove that Gödel's embedding preserves evaluation, 
but a less formal proof is given by \cite{espritosanto_et_al:LIPIcs:2019:10525}

We now give a brief overview of the chapters in this thesis.
\Cref{chap:lambda-calculus} introduces the \lc, with substitution, \br and evaluation relations.
The next chapter, \cref{chap:box-lambda-calculus}, introduces \lab, with substitution, \br and its evaluation relations.
We then discuss the two embeddings in \cref{chap:embeddings}, where we also provide the propositions about the embeddings.
\Cref{chap:agda} explains the basic Agda concepts, the challenges we faced while formalising \lab and their embeddings, and how we solved them.
Finally, \cref{conclusions} concludes the thesis and \cref{relatedwork} gives an overview of related work in the field.
